\chapter[Problemy odwrotne i sterowanie PDE]{Problemy odwrotne i sterowanie optymalne dla równań różniczkowych cząstkowych}
\label{chap:pde32}

\section{Problem prosty, odwrotny i sens stabilności}
Zagadnienie proste polega na wyznaczeniu rozwiązania $y$ PDE
z danych, np. źródła $f$ i współczynników operatora.
W zagadnieniu odwrotnym obserwujemy częściowo $y$ i próbujemy
odzyskać nieznane $f$, współczynniki albo warunki początkowe.
Problem sterowania polega natomiast na doborze danych wejściowych
$u$, aby wymusić pożądaną własność rozwiązania.

\begin{definicja}[Stabilność zagadnienia odwrotnego]
Niech $X,Y$ będą przestrzeniami unormowanymi, $K:X\to Y$
odwzorowaniem obserwacji (niekoniecznie liniowym),
$M\subset X$ klasą dopuszczalnych danych, a $y=Kf$
danymi bez szumu dla $f\in M$.
Zagadnienie odtworzenia $f$ nazywamy stabilnym
w sensie Lipschitza w tych normach, jeśli istnieje
stała $C>0$ spełniająca
$\|f_1-f_2\|_X\le C\|Kf_1-Kf_2\|_Y$ dla wszystkich
$f_1,f_2\in M$.
Warunek ten implikuje różnowartościowość $K|_M$,
a więc istnienie jednoznacznej odwrotności na $K(M)$.
Sama różnowartościowość nie gwarantuje stabilności.
Jeżeli $K$ jest liniowy i $M=X$, jednoznaczność
jest równoważna warunkowi $\ker K=\{0\}$.
Ciągłość odwrotności jest warunkiem słabszym niż
stabilność w sensie Lipschitza; należy odróżniać
te pojęcia, zwłaszcza dla problemów nieliniowych.
\end{definicja}

\begin{przyklad}[Odtwarzanie źródła równania eliptycznego]
\label{ex:pde32-elliptic-inverse}
Niech $\Omega\subset\R^d$ będzie ograniczonym obszarem
Lipschitzowskim, a $A=-\Delta+I$ z jednorodnym warunkiem
Dirichleta. Dla $f\in L^2(\Omega)$ niech $y=Kf=A^{-1}f$.
Twierdzenie Laxa--Milgrama zastosowane do formy
\[
 b(y,v)=\int_\Omega(\nabla y\cdot\nabla v+yv)
 \,\mathrm d\lambda_d(x)
\]
na rzeczywistej przestrzeni $H_0^1(\Omega)$ daje
istnienie i jednoznaczność rozwiązania słabego.
Testując równanie przez $v=y$, otrzymujemy
\[
 \|\nabla y\|_2^2+\|y\|_2^2
 =\int_\Omega fy\,\mathrm d\lambda_d(x)
 \le\|f\|_2\|y\|_2.
\]
Stąd $\|y\|_{H^1}\le\|f\|_2$ (dla $y=0$
nierówność zachodzi bez dzielenia).
Operator $K:L^2(\Omega)\to L^2(\Omega)$ jest zwarty,
ponieważ odwzorowuje zbiory ograniczone w $L^2$
na zbiory ograniczone w $H_0^1(\Omega)$, a zanurzenie
$H_0^1\hookrightarrow L^2$ jest zwarte.
Jeśli $e_k$ jest znormalizowaną funkcją własną laplasjanu
Dirichleta z wartością własną $\lambda_k\to\infty$, to
$Ke_k=(1+\lambda_k)^{-1}e_k$. Mamy
$\|e_k\|_2=1$ i $\|Ke_k\|_2\to0$.
Operator jest różnowartościowy: jeśli $Kf=0$,
to równanie słabe daje $\int_\Omega fv\,\mathrm d\lambda_d(x)=0$
dla wszystkich $v\in H_0^1(\Omega)$, skąd $f=0$
z gęstości $C_c^\infty(\Omega)$ w $L^2(\Omega)$.
Odwrotność $K^{-1}$ na obrazie $K$ nie jest jednak
ciągła nawet w zerze, ponieważ $Ke_k\to0$,
lecz $K^{-1}(Ke_k)=e_k$ nie zbiega do
$K^{-1}(0)=0$.
\end{przyklad}

\begin{uwaga}[Odtwarzanie warunku początkowego równania ciepła]
Dla operatora $-\Delta_D$ obserwacja rozwiązania po czasie
$T>0$ jest dana przez $K_T=\exp(T\Delta_D)$.
Na funkcjach własnych $K_Te_k=e^{-T\lambda_k}e_k$.
Współczynniki wysokoczęstotliwościowe są tłumione
wykładniczo; ich odzyskanie wzmacnia nawet niewielki szum.
Jest to przykład \emph{silnej niestabilności} zagadnienia odwrotnego.
\end{uwaga}


\begin{stwierdzenie}[Zwartość obserwacji cieplnej i brak stabilnego odwrócenia]
\label{prop:pde-v44-heat-inverse}
Niech $\Omega\subset\R^d$ będzie niepustym ograniczonym
obszarem Lipschitzowskim, $H=L^2(\Omega)$ i $T>0$.
Dla laplasjanu Dirichleta niech
$K_T=e^{T\Delta_D}:H\to H$. Wtedy $K_T$ jest
zwartym, samosprzężonym, dodatnim i różnowartościowym
operatorem. Jego odwrotność $K_T^{-1}:\operatorname{Ran}(K_T)\to H$
nie jest ograniczona, przy czym $\operatorname{Ran}(K_T)$
jest gęsty, lecz nie jest domknięty w $H$.
\end{stwierdzenie}
\begin{proof}
Z teorii spektralnej w rozdziale~\ref{chap:pde18}
istnieje baza ortonormalna $(e_k)$ przestrzeni $H$
i wartości własne $0<\lambda_k\to\infty$ operatora
$-\Delta_D$. Mamy $K_Te_k=e^{-T\lambda_k}e_k$.
Obcięcia skończonego rzędu
$K_T^{(N)}f=\sum_{k=1}^Ne^{-T\lambda_k}(f,e_k)e_k$
spełniają
$\|K_T-K_T^{(N)}\|=\sup_{k>N}e^{-T\lambda_k}\to0$.
Istotnie, tożsamość Parsevala daje ograniczenie normy
przez ten kres górny, a testowanie na $e_k$ dla $k>N$
daje nierówność przeciwną. Zatem $K_T$ jest zwarty. Dodatnie rzeczywiste
mnożniki dają samosprzężoność i dodatniość;
ponieważ żaden nie znika, jądro jest zerowe.
Dokładniej, dla $g=\sum_k g_ke_k\in H$ zachodzi
\[
 g\in\operatorname{Ran}(K_T)
 \quad\Longleftrightarrow\quad
 \sum_k e^{2T\lambda_k}|g_k|^2<\infty.
\]
Jeżeli $g=K_Tf$, to $g_k=e^{-T\lambda_k}f_k$,
a warunek wynika z tożsamości Parsevala.
Odwrotnie, gdy szereg jest zbieżny, funkcja
$f=\sum_k e^{T\lambda_k}g_ke_k$ należy do $H$
i spełnia $K_Tf=g$.
W szczególności obraz zawiera wszystkie skończone
kombinacje funkcji własnych, jest więc gęsty.
Nie jest jednak domknięty: wybierając podciąg
$\lambda_{k_j}$ taki, że $e^{T\lambda_{k_j}}\ge j$,
oraz $g_{k_j}=1/j$, a pozostałe współczynniki zerowe,
otrzymujemy $g\in H$ niespełniające warunku
przynależności do obrazu.
Gdyby odwrotność była ograniczona, mielibyśmy
$1=\|e_k\|_H\le C\|K_Te_k\|_H
=C e^{-T\lambda_k}$ dla wszystkich $k$, sprzeczność.
\end{proof}
\begin{uwaga}[Wygładzanie a identyfikowalność]
Zwartość operatora obserwacji wiąże się z wygładzaniem
parabolicznym z rozdziału~\ref{chap:pde20}; stwierdzenie
\ref{prop:pde-v44-heat-wavefront} pokazuje analogiczny
mechanizm z perspektywy zbioru falowego.
Jednoznaczność danych początkowych nie zapewnia ich
stabilnej rekonstrukcji. Twierdzenie~\ref{thm:pde32-tikhonov}
dodaje regularyzację, nie odtwarza utraconej
stabilności nieuregulowanej odwrotności.
\end{uwaga}


\section{Regularyzacja Tichonowa: istnienie i równanie normalne}
\begin{twierdzenie}[Regularyzacja kwadratowa]\label{thm:pde32-tikhonov}
Niech $X,Y$ będą rzeczywistymi przestrzeniami Hilberta,
$K\in\mathcal L(X,Y)$, $y^\delta\in Y$ oraz $\alpha>0$.
Funkcjonał
\begin{equation}\label{eq:pde32-tikhonov}
 J_\alpha(f)=\frac12\|Kf-y^\delta\|_Y^2
 +\frac\alpha2\|f\|_X^2
\end{equation}
ma dokładnie jeden minimizator $f_\alpha^\delta\in X$.
Jest on jedynym rozwiązaniem równania normalnego
\begin{equation}\label{eq:pde32-normal}
 (K^*K+\alpha I)f_\alpha^\delta=K^*y^\delta.
\end{equation}
Dla ustalonego $\alpha>0$ odwzorowanie danych
w minimizator jest liniowe i stabilne w sensie
Lipschitza, ze stałą nie większą niż $1/(2\sqrt\alpha)$.
Jeżeli $\|y^\delta-y\|_Y\le\delta$ i $y=Kf^\dagger$,
gdzie $f^\dagger$ jest rozwiązaniem równania $Kf=y$
o najmniejszej normie, to przy $\alpha\downarrow0$
i $\delta/\sqrt\alpha\to0$ zachodzi
$\|f_\alpha^\delta-f^\dagger\|_X\to0$.
\end{twierdzenie}
\begin{proof}
Funkcjonał jest ściśle wypukły, ciągły i koercywny.
Każdy ciąg minimalizujący jest ograniczony; z refleksywności
przestrzeni Hilberta wybieramy słabo zbieżny podciąg.
Normy po zastosowaniu ograniczonego operatora $K$ są słabo
dolnie półciągłe, więc granica minimalizuje $J_\alpha$.
Ścisła wypukłość składnika $\alpha\|f\|^2/2$ daje
jednoznaczność. Dla dowolnego $f,h\in X$ rozwinięcie kwadratów daje
\[
 \begin{aligned}
 J_\alpha(f+h)-J_\alpha(f)
 &= (Kf-y^\delta,Kh)_Y+\alpha(f,h)_X\\
 &\quad+\tfrac12\|Kh\|_Y^2+\tfrac\alpha2\|h\|_X^2.
 \end{aligned}
\]
Reszta jest ograniczona przez
$(\|K\|^2+\alpha)\|h\|_X^2/2$, więc część liniowa
jest pochodną Frécheta. W minimum pochodna
znika dla każdego $h$, co jest równoważne
~\eqref{eq:pde32-normal}. Odwrotnie, jeśli
równanie normalne zachodzi, to powyższa różnica
jest nieujemna dla każdego $h$, zatem $f$
jest minimizatorem.
Ponadto operator $K^*K+\alpha I$ jest odwracalny
na $X$: forma $(Kf,Kh)_Y+\alpha(f,h)_X$
jest ograniczona i koercywna ze stałą $\alpha$.
Twierdzenie Laxa--Milgrama daje dla dowolnego
$g\in X$ dokładnie jedno rozwiązanie równania
$(K^*K+\alpha I)f=g$. Testując przez $f$,
otrzymujemy $\|f\|_X\le\alpha^{-1}\|g\|_X$,
co dowodzi ograniczoności odwrotności.

Niech $A_0=K^*K\ge0$ i $f_\alpha^0=(A_0+\alpha I)^{-1}A_0f^\dagger$.
Z równania normalnego
\[
 f_\alpha^\delta-f^\dagger
 =-\alpha(A_0+\alpha I)^{-1}f^\dagger
 +(A_0+\alpha I)^{-1}K^*(y^\delta-y).
\]
Ponieważ $(A_0 f,f)_X=\|Kf\|_Y^2$, mamy
$\ker A_0=\ker K$. Zbiór rozwiązań równania $Kf=y$
jest domkniętą przestrzenią afiniczną $f^\dagger+\ker K$,
a rozwiązanie o najmniejszej normie jest jej rzutem
metrycznym z zera. W szczególności
$f^\dagger\perp\ker A_0$.
Korzystamy z twierdzenia spektralnego dla ograniczonego
operatora samosprzężonego (nie zakładamy zwartości $K$).
W przypadku rzeczywistej przestrzeni Hilberta można
stosować rachunek spektralny na jej zespoleniu:
rzeczywiste funkcje borelowskie operatora zachowują
podprzestrzeń rzeczywistą.
Niech $E$ będzie miarą spektralną operatora $A_0$
i $\mu(B)=(E(B)f^\dagger,f^\dagger)_X$ dla zbiorów
borelowskich $B$. Wówczas $\mu(\{0\})=0$, ponieważ
$E(\{0\})$ jest rzutem na $\ker A_0$.
Z rachunku funkcyjnego otrzymujemy
\[
 \|\alpha(A_0+\alpha I)^{-1}f^\dagger\|_X^2
 =\int_{[0,\|A_0\|]}
 \left(\frac{\alpha}{t+\alpha}\right)^2\,\mathrm d\mu(t).
\]
Miara $\mu$ jest skończona, funkcja podcałkowa jest
ograniczona przez $1$ i zbiega do zera $\mu$-prawie
wszędzie. Twierdzenie o zbieżności zmajoryzowanej
daje silną zbieżność pierwszego składnika do zera.
Dla drugiego składnika połóżmy
$B_\alpha=(A_0+\alpha I)^{-1}K^*$.
Ponieważ $(A_0+\alpha I)^{-1}$ jest samosprzężony,
\[
 B_\alpha B_\alpha^*
 =(A_0+\alpha I)^{-1}A_0(A_0+\alpha I)^{-1}
 =A_0(A_0+\alpha I)^{-2}.
\]
Z tożsamości $\|B_\alpha\|^2=\|B_\alpha B_\alpha^*\|$
i rachunku spektralnego wynika
\[
 \|B_\alpha\|^2
 \le\sup_{t\ge0}\frac{t}{(t+\alpha)^2}
 =\frac1{4\alpha}.
\]
Ostatnią równość daje nierówność
$(t+\alpha)^2\ge4t\alpha$, z równością dla $t=\alpha$.
Dla dowolnych danych $y_1,y_2\in Y$ ich minimizatory
$f_{\alpha,1},f_{\alpha,2}$ spełniają
$f_{\alpha,1}-f_{\alpha,2}=B_\alpha(y_1-y_2)$,
więc
\[
 \|f_{\alpha,1}-f_{\alpha,2}\|_X
 \le\frac{\|y_1-y_2\|_Y}{2\sqrt\alpha}.
\]
W szczególności błąd od szumu nie przekracza
$\delta/(2\sqrt\alpha)$. Razem z zanikiem
składnika obciążenia przy $\alpha\downarrow0$
oraz założeniem $\delta/\sqrt\alpha\to0$
daje to zbieżność z tezy.
\end{proof}

\begin{stwierdzenie}[Oszacowanie ze źródłowym warunkiem regularności]
\label{prop:pde32-source}
Przy założeniach twierdzenia~\ref{thm:pde32-tikhonov}, jeśli
$f^\dagger=K^*w$ dla pewnego $w\in Y$, to
\begin{equation}\label{eq:pde32-rate}
 \|f_\alpha^\delta-f^\dagger\|_X
 \le \frac{\sqrt\alpha}{2}\|w\|_Y
 +\frac\delta{2\sqrt\alpha}.
\end{equation}
Wybór $\alpha\asymp\delta$ daje błąd $O(\sqrt\delta)$,
przy ustalonym $\|w\|_Y$.
\end{stwierdzenie}
\begin{proof}
Z rozkładu błędu w dowodzie twierdzenia~\ref{thm:pde32-tikhonov} i warunku $f^\dagger=K^*w$
wynika, że składnik obciążenia jest równy
$-\alpha(A_0+\alpha I)^{-1}K^*w=-\alpha B_\alpha w$.
Wykazane tam oszacowanie $\|B_\alpha\|\le1/(2\sqrt\alpha)$
daje
\[
 \|\alpha B_\alpha w\|_X
 \le\frac{\sqrt\alpha}{2}\|w\|_Y .
\]
Część odpowiadającą szumowi oszacowano w poprzednim dowodzie.
\end{proof}


\begin{twierdzenie}[Ułamkowy warunek źródłowy i tempo Tichonowa]
\label{thm:pde32-fractional-source}
Niech $X,Y$ będą rzeczywistymi przestrzeniami Hilberta,
$K\in\mathcal L(X,Y)$, $A_0=K^*K$,
$0<\nu\le1$, $w\in X$ oraz
$f^\dagger=A_0^\nu w$. Przyjmijmy
$y=Kf^\dagger$ i $\|y^\delta-y\|_Y\le\delta$.
Wtedy $f^\dagger$ jest rozwiązaniem $Kf=y$
o najmniejszej normie, a minimizator
\eqref{eq:pde32-tikhonov} spełnia
\begin{equation}\label{eq:pde32-fractional-rate}
 \|f_\alpha^\delta-f^\dagger\|_X
 \le C_\nu\alpha^\nu\|w\|_X+
 \frac{\delta}{2\sqrt\alpha},\qquad
 C_\nu=\sup_{r\ge0}\frac{r^\nu}{1+r}<\infty.
\end{equation}
Jeżeli $\delta\downarrow0$ i $\|w\|_X$ pozostaje
ograniczona, to wybór
$\alpha\asymp\delta^{2/(2\nu+1)}$
daje błąd $O(\delta^{2\nu/(2\nu+1)})$.
\end{twierdzenie}
\begin{proof}
Z twierdzenia spektralnego dla dodatniego ograniczonego
operatora samosprzężonego $A_0$ definiujemy
$A_0^\nu$. Dla $\nu>0$ i $v\in\ker A_0=\ker K$
rachunek funkcyjny daje $A_0^\nu v=0$.
Ponieważ $A_0^\nu$ jest samosprzężony,
$(f^\dagger,v)_X=(w,A_0^\nu v)_X=0$.
Każde rozwiązanie równania $Kf=y$ ma postać
$f=f^\dagger+v$, gdzie $v\in\ker K$, a wtedy
$\|f\|_X^2=\|f^\dagger\|_X^2+\|v\|_X^2$.
Stąd $f^\dagger$ jest jedynym rozwiązaniem
o najmniejszej normie.
Rozkład błędu w dowodzie
\ref{thm:pde32-tikhonov} ma postać
\[
 f_\alpha^\delta-f^\dagger
 =-\alpha(A_0+\alpha I)^{-1}A_0^\nu w
 +(A_0+\alpha I)^{-1}K^*(y^\delta-y).
\]
Twierdzenie spektralne daje
\[
 \|\alpha(A_0+\alpha I)^{-1}A_0^\nu\|
 \le\sup_{t\ge0}\frac{\alpha t^\nu}{t+\alpha}
 =\alpha^\nu\sup_{r\ge0}\frac{r^\nu}{1+r}
 =C_\nu\alpha^\nu.
\]
Dla $0<\nu<1$ pochodna funkcji
$g_\nu(r)=r^\nu/(1+r)$ dla $r>0$ wynosi
\[
 g_\nu'(r)=
 \frac{r^{\nu-1}(\nu-(1-\nu)r)}{(1+r)^2}.
\]
Funkcja rośnie do $r=\nu/(1-\nu)$, a następnie
maleje; ponadto dąży do zera przy $r\downarrow0$
oraz przy $r\to\infty$.
Podstawienie punktu maksimum daje
\[
 C_\nu=\nu^\nu(1-\nu)^{1-\nu}
 \quad(0<\nu<1).
\]
Dla $\nu=1$ funkcja $r/(1+r)$ rośnie do $1$,
zatem $C_1=1$.
Część pochodząca od szumu ma normę nie większą
niż $\delta/(2\sqrt\alpha)$, jak wykazano w
\ref{thm:pde32-tikhonov}. To dowodzi oszacowania.
Aby oba człony oszacowania miały ten sam rząd,
przyjmujemy $\alpha^{\nu+1/2}\asymp\delta$.
Dokładniej, jeśli dla dostatecznie małych $\delta>0$
zachodzi $c_1\delta^{2/(2\nu+1)}\le\alpha
\le c_2\delta^{2/(2\nu+1)}$ z $c_1,c_2>0$, to
\[
 \alpha^\nu\le c_2^\nu\delta^{2\nu/(2\nu+1)},
 \qquad
 \frac{\delta}{\sqrt\alpha}\le
 c_1^{-1/2}\delta^{2\nu/(2\nu+1)}.
\]
To uzasadnia deklarowane tempo przy jednostajnym
ograniczeniu normy $w$. Nie stwierdzamy, że ten wybór
parametru minimalizuje rzeczywisty błąd dla każdych danych.
\end{proof}

\begin{uwaga}[Zakres tempa]
Jest to oszacowanie warunkowe: bez założenia
$f^\dagger\in\operatorname{Ran}(A_0^\nu)$
nie można przypisać wszystkim danym tempa
$\delta^{2\nu/(2\nu+1)}$. Ponadto w standardowej
regularyzacji Tichonowa wykładnik $\nu>1$ nie
prowadzi do analogicznego polepszenia rzędu
członu obciążenia. Istotnie, dla
$f^\dagger=A_0^\nu w$ i danych dokładnych
$y=Kf^\dagger$ mamy
\[
 \|f_\alpha^0-f^\dagger\|_X
 \le\alpha\|A_0\|^{\nu-1}\|w\|_X,
\]
gdyż $t^\nu/(t+\alpha)\le t^{\nu-1}$
dla $t>0$. Co więcej, z rachunku spektralnego
i twierdzenia o zbieżności zmajoryzowanej wynika
\[
 \frac{f_\alpha^0-f^\dagger}{\alpha}
 =-(A_0+\alpha I)^{-1}A_0^\nu w
 \longrightarrow -A_0^{\nu-1}w
 \quad\text{w }X .
\]
Rzeczywiście, mnożniki $t^\nu/(t+\alpha)$
zbiegają punktowo na widmie do $t^{\nu-1}$
(również w zerze) i są jednostajnie ograniczone
przez $\|A_0\|^{\nu-1}$.
Jeżeli $f^\dagger\ne0$, to $A_0^{\nu-1}w\ne0$,
ponieważ $A_0 A_0^{\nu-1}w=f^\dagger$.
Wówczas błąd dla danych dokładnych jest dokładnie
rzędu $\alpha$, nie zaś $o(\alpha)$.
Jest to nasycenie rzędu obciążenia; nie jest to
jeszcze twierdzenie o optymalnym tempie rekonstrukcji
dla dowolnej reguły doboru parametru i klasy szumów.
\end{uwaga}

\begin{przyklad}[Zbyt mała regularyzacja]
Dla $K_T=\exp(T\Delta_D)$ mamy
$\sigma_k=e^{-T\lambda_k}$. Dla dokładnego stanu początkowego $f^\dagger=e_k$
i obserwacji $y^\delta=(\sigma_k+\delta)e_k$
nieuregulowana odwrotność daje
$(1+\delta/\sigma_k)e_k$, zatem mnoży szum przez
$1/\sigma_k=e^{T\lambda_k}$.
Z równania normalnego i samosprzężoności $K_T$
wynika natomiast
\[
 f_\alpha^\delta
 =\frac{\sigma_k(\sigma_k+\delta)}
 {\sigma_k^2+\alpha}e_k,\qquad
 f_\alpha^\delta-f^\dagger
 =\frac{-\alpha+\sigma_k\delta}
 {\sigma_k^2+\alpha}e_k .
\]
Mnożnik szumu $\sigma_k/(\sigma_k^2+\alpha)$
jest ograniczony przez $1/(2\sqrt\alpha)$.
Dla ustalonej $\alpha>0$ i $\delta=0$ błąd ma normę
$\alpha/(\sigma_k^2+\alpha)$, która dąży do $1$
przy $k\to\infty$. Regularyzacja tłumi więc
również część sygnału, nie tylko szum. W celu uzyskania zbieżności do rozwiązania dokładnego przy
zanikającym szumie trzeba dobierać $\alpha=\alpha(\delta)$;
parametr stały ogranicza wpływ szumu, ale dla
niezerowego dokładnego rozwiązania o najmniejszej
normie pozostawia niezanikający błąd regularyzacji:
$\alpha(A_0+\alpha I)^{-1}f^\dagger\ne0$
dla $\alpha>0$, ponieważ $(A_0+\alpha I)^{-1}$
jest różnowartościowy.
\end{przyklad}

\section{Sterowanie rozłożone w obszarze: operator stanu}
Niech $\Omega\subset\R^d$ będzie ograniczonym obszarem
Lipschitzowskim, $V=H_0^1(\Omega)$ i $H=L^2(\Omega)$
nad ciałem liczb rzeczywistych. W dalszej części
przyjmujemy normę $\|v\|_V=\|\nabla v\|_2$,
równoważną standardowej normie $H^1$ na $V$
na mocy nierówności Poincarégo
$\|v\|_H\le C_P\|v\|_V$.
Normę na przestrzeni funkcjonałów liniowych i ciągłych
$V'=H^{-1}(\Omega)$
wyznaczamy względem tej normy na $V$.
Przyjmujemy mierzalną symetryczną macierz
$A\in L^\infty(\Omega;\R^{d\times d})$
i $\xi^{\mathsf T}A(x)\xi\ge\lambda|\xi|^2$ dla
każdego $\xi$ i prawie każdego $x$, z $\lambda>0$.
Niech $c\in L^\infty(\Omega)$, $c\ge0$.
Forma
$\mathfrak a(y,v)=\int_\Omega
(A\nabla y\cdot\nabla v+cyv)\,\mathrm d\lambda_d(x)$
jest symetryczna, ograniczona i koercywna na $V$.
Dokładniej, z nierówności Cauchy'ego--Schwarza
i nierówności Poincarégo mamy
\[
 \begin{aligned}
 |\mathfrak a(y,v)|
 &\le(\|A\|_{L^\infty}+C_P^2\|c\|_\infty)
 \|y\|_V\|v\|_V,\\
 \mathfrak a(v,v)&\ge\lambda\|v\|_V^2,
 \end{aligned}
\]
gdzie normę macierzy w $L^\infty$ rozumiemy
jako istotny kres górny normy operatorowej macierzy.
Niech $f\in V'$ oraz $u\in H$. Funkcjonał
$v\mapsto\langle f,v\rangle_{V',V}+(u,v)_H$
jest liniowy i ograniczony przez
$(\|f\|_{V'}+C_P\|u\|_H)\|v\|_V$.
Z twierdzenia Laxa--Milgrama istnieje zatem
dokładnie jedno $y(u)\in V$ spełniające równanie
\begin{equation}\label{eq:pde32-state}
 \mathfrak a(y(u),v)=\langle f,v\rangle_{V',V}
 +\int_\Omega uv\,\mathrm d\lambda_d(x)
 \quad\text{dla każdego }v\in V.
\end{equation}
Testując równanie przez $v=y(u)$, otrzymujemy
\[
 \|y(u)\|_V\le\lambda^{-1}
 (\|f\|_{V'}+C_P\|u\|_H).
\]
Kładziemy $S(u)=y(u)-y(0)$. Odejmując równania
dla $y(u)$ oraz $y(0)$, mamy
$\mathfrak a(S(u),v)=(u,v)_H$ dla każdego $v\in V$.
Liniowość prawej strony oraz jednoznaczność
rozwiązania dają liniowość $S$.
Z koercywności, po podstawieniu $v=S(u)$,
wynika $\|S(u)\|_V\le C_P\lambda^{-1}\|u\|_H$.
W szczególności $u\mapsto y(u)=y(0)+S(u)$
jest afiniczne i ciągłe jako odwzorowanie $H\to V$.
Traktowany jako operator $H\to H$, operator $S$
jest także zwarty dzięki zwartemu zanurzeniu $V\hookrightarrow H$.

\begin{twierdzenie}[Istnienie sterowania optymalnego]\label{thm:pde32-control}
Niech $y_d\in H$, $\alpha>0$ i $U_{\mathrm{ad}}\subset H$
będzie niepustym, domkniętym i wypukłym podzbiorem.
Wtedy problem
\begin{equation}\label{eq:pde32-opt}
 \min_{u\in U_{\mathrm{ad}}}
 \mathcal J(u),\qquad
 \mathcal J(u)=\frac12\|y(u)-y_d\|_H^2
 +\frac\alpha2\|u\|_H^2
\end{equation}
ma dokładnie jedno rozwiązanie optymalne $\bar u$.
\end{twierdzenie}
\begin{proof}
Dla dowolnego ciągu minimalizującego $(u_n)$ nierówność
$\alpha\|u_n\|_H^2/2\le\mathcal J(u_n)$ daje ograniczoność.
Istnieje słabo zbieżny podciąg $u_n\rightharpoonup\bar u$ w $H$.
Każdy domknięty wypukły podzbiór przestrzeni Hilberta jest
słabo domknięty, więc $\bar u\in U_{\mathrm{ad}}$.
Ciągłość liniowa $S$ daje słabą zbieżność
$y(u_n)\rightharpoonup y(\bar u)$ w $H$.
Słaba dolna półciągłość norm daje
$\mathcal J(\bar u)\le\liminf\mathcal J(u_n)$.
Dla $u,v\in H$ i $0\le\theta\le1$ z afiniczności
stanu i tożsamości dla kwadratu normy w przestrzeni
Hilberta otrzymujemy
\[
 \begin{aligned}
 \mathcal J((1-\theta)u+\theta v)
 &=(1-\theta)\mathcal J(u)+\theta\mathcal J(v)\\
 &\quad-\frac{\theta(1-\theta)}2
 \bigl(\|S(u-v)\|_H^2+\alpha\|u-v\|_H^2\bigr).
 \end{aligned}
\]
Dla $u\ne v$ i $0<\theta<1$ ostatni człon jest
ściśle ujemny, co daje ścisłą wypukłość.
Dwa różne minimizatory w $U_{\mathrm{ad}}$
miałyby dopuszczalny środek o mniejszej wartości
funkcjonału, co jest sprzecznością.
Ponadto każdy ciąg minimalizujący zbiega silnie
do $\bar u$: podstawiając $u=\bar u$, $v=u_n$,
$\theta=1/2$ i korzystając z minimalności
$\bar u$, otrzymujemy
\[
 \frac{\alpha}{8}\|u_n-\bar u\|_H^2
 \le \frac12\bigl(\mathcal J(u_n)-\mathcal J(\bar u)\bigr)
 \longrightarrow0.
\]
\end{proof}

\section{Równanie sprzężone i warunki optymalności}
\begin{twierdzenie}[Warunek pierwszego rzędu i rzut na sterowania dopuszczalne]
\label{thm:pde32-adjoint}
Niech $\bar u\in U_{\mathrm{ad}}$ będzie dowolnym
sterowaniem dopuszczalnym, $\bar y=y(\bar u)$,
a $\bar p\in V$ niech będzie jedynym rozwiązaniem
równania sprzężonego
\begin{equation}\label{eq:pde32-adj}
 \mathfrak a(v,\bar p)
 =\int_\Omega(\bar y-y_d)v\,\mathrm d\lambda_d(x)
 \qquad(v\in V).
\end{equation}
Wówczas $\bar u$ jest rozwiązaniem
~\eqref{eq:pde32-opt} wtedy i tylko wtedy, gdy
spełnia nierówność wariacyjną
\begin{equation}\label{eq:pde32-variational-control}
 \int_\Omega(\bar p+\alpha\bar u)(v-\bar u)
 \,\mathrm d\lambda_d(x)\ge0
 \quad\text{dla każdego }v\in U_{\mathrm{ad}},
\end{equation}
a równoważnie
$\bar u=P_{U_{\mathrm{ad}}}(-\bar p/\alpha)$,
gdzie $P_{U_{\mathrm{ad}}}$ jest rzutem metrycznym na
niepusty domknięty wypukły podzbiór przestrzeni Hilberta.
Jeżeli $U_{\mathrm{ad}}=\{v\in L^2(\Omega):a\le v\le b
\text{ prawie wszędzie}\}$ dla liczb $a<b$,
to powyższe warunki są równoważne warunkowi
$\bar u(x)=\min\{b,\max\{a,-\bar p(x)/\alpha\}\}$
dla prawie każdego $x\in\Omega$.
\end{twierdzenie}
\begin{proof}
Forma $\mathfrak a$ jest ograniczona i koercywna,
a funkcjonał $v\mapsto(\bar y-y_d,v)_H$ jest
ograniczony na $V$ przez nierówność Poincarégo.
Twierdzenie Laxa--Milgrama zastosowane do formy
$(p,v)\mapsto\mathfrak a(v,p)$ daje istnienie
i jednoznaczność $\bar p$.
Ponieważ $y(\bar u+h)=\bar y+S h$, pochodna
Frécheta stanu w kierunku $h\in H$ jest równa
$z=S h$ i spełnia
$\mathfrak a(z,v)=\int_\Omega hv\,\mathrm d\lambda_d(x)$.
Rozwijając kwadraty norm w funkcjonale, otrzymujemy
\[
 \begin{aligned}
 \mathcal J(\bar u+h)-\mathcal J(\bar u)
 &= (\bar y-y_d,S h)_H+\alpha(\bar u,h)_H\\
 &\quad+\tfrac12\|S h\|_H^2+\tfrac\alpha2\|h\|_H^2.
 \end{aligned}
\]
Pierwsze dwa składniki wyznaczają ograniczony
funkcjonał liniowy względem $h$, a suma pozostałych
jest $O(\|h\|_H^2)$ z ograniczoności $S$.
Zatem $\mathcal J'(\bar u)h=(\bar y-y_d,z)_H+
\alpha(\bar u,h)_H$. Z równań dla $z$ i $\bar p$
wynika
$(\bar y-y_d,z)_H=\mathfrak a(z,\bar p)=
(h,\bar p)_H$. Stąd $\mathcal J'(\bar u)h=
(\bar p+\alpha\bar u,h)_H$.
Załóżmy teraz, że $\bar u$ jest minimizatorem.
Dla $v\in U_{\mathrm{ad}}$ z wypukłości zbioru
dopuszczalnego wynika
$\bar u+t(v-\bar u)\in U_{\mathrm{ad}}$ dla $t\in[0,1]$.
Funkcja $t\mapsto\mathcal J(\bar u+t(v-\bar u))$
ma zatem minimum w $t=0$ na $[0,1]$.
Jej pochodna prawostronna w zerze jest nieujemna
i równa $\mathcal J'(\bar u)(v-\bar u)$,
co daje~\eqref{eq:pde32-variational-control}.
Odwrotnie, jeśli nierówność wariacyjna zachodzi,
to dla $h=v-\bar u$, $v\in U_{\mathrm{ad}}$,
powyższe rozwinięcie funkcjonału daje
\[
 \mathcal J(v)-\mathcal J(\bar u)
 \ge\tfrac\alpha2\|v-\bar u\|_H^2\ge0.
\]
Dowodzi to optymalności i jednoznaczności.
Dla niepustego domkniętego wypukłego zbioru
$C\subset H$ rzut metryczny $z=P_Cw$ istnieje
i jest jednoznaczny z twierdzenia o rzucie
w przestrzeni Hilberta. Jego charakterystyka
\[
 z=P_Cw\quad\Longleftrightarrow\quad
 z\in C,\quad (z-w,v-z)_H\ge0\quad(v\in C)
\]
wynika w jedną stronę z różniczkowania
$\|z+t(v-z)-w\|_H^2$ w $t=0$ od prawej,
a w drugą z tożsamości
$\|v-w\|_H^2=\|z-w\|_H^2+
2(z-w,v-z)_H+\|v-z\|_H^2$.
Podstawiając $C=U_{\mathrm{ad}}$ i $w=-\bar p/\alpha$
oraz dzieląc nierówność wariacyjną przez
$\alpha>0$, otrzymujemy wzór na rzut.

W przypadku ograniczeń punktowych zbiór
$U_{\mathrm{ad}}$ jest niepusty (zawiera funkcję
stałą $a$), wypukły i domknięty w $L^2$:
z ciągu zbieżnego w $L^2$ można wybrać
podciąg zbieżny prawie wszędzie, więc ograniczenia
$a\le v\le b$ przechodzą na granicę.
Dla $w\in H$ określamy
$q(x)=\min\{b,\max\{a,w(x)\}\}$.
Funkcja $q$ jest mierzalna i ograniczona, należy
do $U_{\mathrm{ad}}$, ponieważ $\lambda_d(\Omega)<\infty$.
Dla każdej liczby $r\in[a,b]$ zachodzi
$(q(x)-w(x))(r-q(x))\ge0$: gdy $w(x)<a$,
oba czynniki są nieujemne, gdy $w(x)>b$,
oba są niedodatnie, a w pozostałym przypadku
pierwszy czynnik jest zerowy.
Podstawiając $r=v(x)$ i całkując, dostajemy
$(q-w,v-q)_H\ge0$ dla wszystkich $v\in U_{\mathrm{ad}}$.
Z charakterystyki rzutu $q=P_{U_{\mathrm{ad}}}w$,
co daje podany wzór punktowy.
\end{proof}

\begin{uwaga}[Dlaczego równanie sprzężone jest użyteczne?]
Bez zmiennej sprzężonej obliczenie gradientu kosztu
wymagałoby rozwiązania nowego PDE dla każdego kierunku $h$.
Równanie~\eqref{eq:pde32-adj} wymaga jednego dodatkowego
rozwiązania PDE i daje gradient jako funkcję $\alpha u+p$.
To podstawowy mechanizm metod numerycznych sterowania.
\end{uwaga}

\section{Sterowanie równaniem ciepła i czas wsteczny}
\begin{stwierdzenie}[Równanie sprzężone dla parabolicznego sterowania rozłożonego]
\label{prop:pde32-heat-adjoint}
Niech $T>0$, $\Omega$ będzie ograniczonym obszarem
Lipschitzowskim, $u_0\in L^2(\Omega)$, $y_d\in
L^2((0,T)\times\Omega)$ oraz $\alpha>0$.
Dla $v\in L^2((0,T)\times\Omega)$ niech $y(v)$ będzie
słabym rozwiązaniem
$y_t-\Delta y=v$ w $(0,T)\times\Omega$,
$y=0$ na $(0,T)\times\partial\Omega$ i $y(0)=u_0$.
Koszt
$\mathcal J_T(v)=\frac12\|y(v)-y_d\|_{L^2((0,T)\times\Omega)}^2
+\frac\alpha2\|v\|_{L^2((0,T)\times\Omega)}^2$
ma pochodną
$\mathcal J_T'(v)h=
\int_0^T\!\int_\Omega(\alpha v+p)h
\,\mathrm d\lambda_d(x)\,\mathrm d\lambda_1(t)$,
gdzie $p$ jest rozwiązaniem równania sprzężonego
\[
 -p_t-\Delta p=y(v)-y_d,
 \qquad p(T)=0,\qquad p|_{\partial\Omega}=0.
\]
\end{stwierdzenie}
\begin{proof}
Piszemy $Q=(0,T)\times\Omega$, $V=H_0^1(\Omega)$,
$H=L^2(\Omega)$ oraz $V'=H^{-1}(\Omega)$.
Z Twierdzenia~\ref{thm:pde20-existence}
w rozdziale~\ref{chap:pde20}
dla każdego $v\in L^2(Q)$ istnieje dokładnie jedno
rozwiązanie słabe w przestrzeni
\[
 W=\{w\in L^2(0,T;V):w_t\in L^2(0,T;V')\},
\]
zanurzonej w $C([0,T];H)$. Wartości początkowe
i końcowe rozumiemy w tym ciągłym reprezentancie.
Liniowość równania i jednoznaczność dają
$y(v+h)-y(v)=z_h$, gdzie $z_h\in W$ spełnia
\[
 \langle(z_h)_t,\phi\rangle_{V',V}
 +(\nabla z_h,\nabla\phi)_{L^2(\Omega)}
 =(h,\phi)_H,\qquad z_h(0)=0.
\]
Równość zachodzi w $V'$ dla prawie każdego $t$.
Testując przez $z_h(t)$ i stosując nierówność
Poincarégo oraz nierówność Younga, otrzymujemy
\[
 \tfrac12\frac{\mathrm d}{\mathrm dt}\|z_h\|_H^2
 +\tfrac12\|z_h\|_V^2
 \le\tfrac{C_P^2}{2}\|h\|_H^2
\]
dla prawie każdego $t$; pochodna normy jest
uzasadniona formułą energii w trójce Gelfanda.
Po całkowaniu i użyciu $z_h(0)=0$ dostajemy
$\|z_h\|_{L^2(0,T;V)}\le C_P\|h\|_{L^2(Q)}$,
a więc $\|z_h\|_{L^2(Q)}\le C_P^2\|h\|_{L^2(Q)}$.
Operator $h\mapsto z_h$ jest liniowy i ograniczony,
stąd jest pochodną Frécheta stanu do $L^2(Q)$.

Aby skonstruować $p$, rozwiązujemy zwykłe
zagadnienie paraboliczne
$q_s-\Delta q=(y(v)-y_d)(T-s)$ z $q(0)=0$
i jednorodnym warunkiem brzegowym.
Prawa strona należy do $L^2(Q)$.
Podstawienie $p(t)=q(T-t)$ daje $p\in W$
spełniające równanie sprzężone z $p(T)=0$;
to samo odwrócenie czasu dowodzi jednoznaczności.
Nie jest to niestabilne odwracanie ewolucji
cieplnej z zadanym stanem obserwowanym,
lecz zwykłe zagadnienie początkowe po zmianie czasu.

Testujemy słabe równanie $z_h$ przez $p(t)$,
a słabe równanie $p$ przez $z_h(t)$ i całkujemy.
Wszystkie parowania są całkowalne z nierówności
Cauchy'ego--Schwarza w czasie. Odejmując równości,
otrzymujemy
\[
 \begin{aligned}
 &\int_0^T\bigl((h,p)_H-(y(v)-y_d,z_h)_H\bigr)
 \,\mathrm d\lambda_1(t)\\
 &\qquad=\int_0^T
 \bigl(\langle(z_h)_t,p\rangle_{V',V}
 +\langle p_t,z_h\rangle_{V',V}\bigr)
 \,\mathrm d\lambda_1(t)\\
 &\qquad=(z_h(T),p(T))_H-(z_h(0),p(0))_H=0.
 \end{aligned}
\]
W ostatnim kroku użyliśmy formuły całkowania
przez części w trójce Gelfanda, wynikającej
z Lematu~\ref{lem:pde20-energy} przez polaryzację:
dla $z_h+p$ oraz $z_h-p$ odejmujemy tożsamości
energii i dzielimy przez dwa, uzyskując tożsamość
dla przyrostu iloczynu skalarnego $(z_h,p)_H$.
Odsyłamy do rozdziału~\ref{chap:pde20}.
Wyrazy przestrzenne znoszą się
dzięki symetrii formy Dirichleta.
Rozwinięcie kwadratów w funkcjonale daje
\[
 \begin{aligned}
 \mathcal J_T(v+h)-\mathcal J_T(v)
 &= (y(v)-y_d,z_h)_{L^2(Q)}
   +\alpha(v,h)_{L^2(Q)}\\
 &\quad+\tfrac12\|z_h\|_{L^2(Q)}^2+
 \tfrac\alpha2\|h\|_{L^2(Q)}^2 .
 \end{aligned}
\]
Reszta jest $O(\|h\|_{L^2(Q)}^2)$ z oszacowania
operatora stanu. Podstawiając wykazaną tożsamość
do części liniowej, otrzymujemy pochodną Frécheta
z tezy.
\end{proof}

\section{Odtwarzanie współczynnika i granice identyfikowalności}
\begin{stwierdzenie}[Pochodna względem współczynnika dyfuzji]
\label{prop:pde32-coeff-gradient}
Niech $\Omega$ będzie ograniczonym obszarem Lipschitzowskim,
$0<\lambda\le a\le\Lambda$ prawie wszędzie,
$a\in L^\infty(\Omega)$, $f\in H^{-1}(\Omega)$.
Niech $y(a)\in H_0^1(\Omega)$ będzie jedynym rozwiązaniem
$-\diver(a\nabla y(a))=f$.
Dla $h\in L^\infty(\Omega)$ i $|\eps|$ dostatecznie małych
niech $a_\eps=a+\eps h$. Wtedy istnieje pochodna kierunkowa
$z=\lim_{\eps\to0}(y(a_\eps)-y(a))/\eps$ w $H_0^1(\Omega)$
i spełnia
\begin{equation}\label{eq:pde32-coeff-linear}
 \int_\Omega a\nabla z\cdot\nabla v
 \,\mathrm d\lambda_d(x)
 =-\int_\Omega h\nabla y(a)\cdot\nabla v
 \,\mathrm d\lambda_d(x)\qquad(v\in H_0^1(\Omega)).
\end{equation}
Jeżeli $y_d\in L^2(\Omega)$,
$\mathcal F(a)=\frac12\|y(a)-y_d\|_2^2$,
a $p\in H_0^1(\Omega)$ jest jedynym rozwiązaniem
$-\diver(a\nabla p)=y(a)-y_d$ w $H^{-1}$,
z jednorodnym warunkiem Dirichleta, to
$\mathcal F'(a)h=-\int_\Omega
h\nabla y(a)\cdot\nabla p\,\mathrm d\lambda_d(x)$.
\end{stwierdzenie}
\begin{proof}
Pracujemy z normami $\|v\|_V=\|\nabla v\|_2$
oraz odpowiadającą jej normą na $V'=H^{-1}(\Omega)$.
Dla $|\eps|\|h\|_\infty\le\lambda/2$ mamy
$a_\eps\ge\lambda/2$ prawie wszędzie.
Istnienie i jednoznaczność $y(a_\eps)$ wynikają
z twierdzenia Laxa--Milgrama.
Testując równanie przez $y(a_\eps)$, dostajemy
$\|y(a_\eps)\|_V\le2\lambda^{-1}\|f\|_{V'}$.
Oznaczmy $d_\eps=y(a_\eps)-y(a)$.
Po odjęciu równań stanów otrzymujemy
\[
 \int_\Omega a\nabla d_\eps\cdot\nabla v
 \,\mathrm d\lambda_d(x)
 =-\eps\int_\Omega h\nabla y(a_\eps)\cdot\nabla v
 \,\mathrm d\lambda_d(x).
\]
Dla $v=d_\eps$ koercywność i nierówność
Cauchy'ego--Schwarza dają
\[
 \|d_\eps\|_V
 \le 2\lambda^{-2}|\eps|\|h\|_\infty\|f\|_{V'}.
\]
Gdy $d_\eps=0$, oszacowanie zachodzi bez dzielenia.
Iloraz różnicowy $z_\eps$ spełnia
$\int_\Omega a\nabla z_\eps\cdot\nabla v\,\mathrm d\lambda_d(x)
=-\int_\Omega h\nabla y(a_\eps)\cdot\nabla v\,\mathrm d\lambda_d(x)$.
Prawa strona równania~\eqref{eq:pde32-coeff-linear}
jest ograniczonym funkcjonałem na $V$, zatem
twierdzenie Laxa--Milgrama daje dokładnie jedno $z$.
Odejmując równanie dla $z$ od równania dla
$z_\eps$ i testując przez $z_\eps-z$, dostajemy
\[
 \|z_\eps-z\|_V
 \le\lambda^{-1}\|h\|_\infty\|d_\eps\|_V
 \le2\lambda^{-3}|\eps|\|h\|_\infty^2\|f\|_{V'}.
\]
W szczególności $z_\eps\to z$ silnie w $V$.
Równanie dla $p$ ma dokładnie jedno rozwiązanie
w $V$ na mocy twierdzenia Laxa--Milgrama, ponieważ
$y(a)-y_d\in L^2(\Omega)\hookrightarrow V'$.
Ponadto
\[
 \frac{\mathcal F(a_\eps)-\mathcal F(a)}{\eps}
 =(y(a)-y_d,z_\eps)_2+\frac{\eps}{2}\|z_\eps\|_2^2.
\]
Z silnej zbieżności $z_\eps$ w $V$, a więc w $L^2$,
wynika istnienie pochodnej kierunkowej kosztu.
Z definicji $p$ i~\eqref{eq:pde32-coeff-linear}
$\mathcal F'(a)h=(y(a)-y_d,z)_2
=\int_\Omega a\nabla p\cdot\nabla z\,\mathrm d\lambda_d(x)
=-\int_\Omega h\nabla y(a)\cdot\nabla p\,\mathrm d\lambda_d(x)$.
\end{proof}


\begin{wniosek}[Różniczkowalność w normie $L^\infty$]
\label{cor:pde32-frechet}
Przy założeniach stwierdzenia~\ref{prop:pde32-coeff-gradient}
operator rozwiązania $a\mapsto y(a)$ jest różniczkowalny
w sensie Frécheta jako odwzorowanie z otwartego zbioru
$\{a\in L^\infty(\Omega):\operatorname*{ess\,inf}a>0\}$
do $H_0^1(\Omega)$. Jego pochodna w kierunku $h\in L^\infty$
jest jedynym rozwiązaniem~\eqref{eq:pde32-coeff-linear}.
W poniższym oszacowaniu normą na $H_0^1(\Omega)$
jest $\|\nabla\cdot\|_2$, a na $H^{-1}(\Omega)$
odpowiadająca jej norma funkcjonałów liniowych i ciągłych.
Dla $\|h\|_\infty\le\lambda/2$ istnieje stała $C$
zależna wyłącznie od $\lambda$, dla której
\[
 \|y(a+h)-y(a)-z_h\|_{H_0^1}
 \le C\|h\|_\infty^2\|f\|_{H^{-1}}.
\]
\end{wniosek}
\begin{proof}
Oznaczmy $y_h=y(a+h)$, $y=y(a)$ i $\delta y=y_h-y$.
Odejmowanie równań stanów daje
$\int_\Omega a\nabla\delta y\cdot\nabla v\,\mathrm d\lambda_d(x)
=-\int_\Omega h\nabla y_h\cdot\nabla v\,\mathrm d\lambda_d(x)$.
Ponieważ $a+h\ge\lambda/2$, z oszacowań energetycznych
wynikają $\|y_h\|_{H_0^1}\le2\lambda^{-1}\|f\|_{H^{-1}}$
i $\|\delta y\|_{H_0^1}\le
2\lambda^{-2}\|h\|_\infty\|f\|_{H^{-1}}$.
Kandydat na wartość pochodnej $z_h$ spełnia
$\int_\Omega a\nabla z_h\cdot\nabla v\,\mathrm d\lambda_d(x)
=-\int_\Omega h\nabla y\cdot\nabla v\,\mathrm d\lambda_d(x)$.
Dla reszty $r_h=\delta y-z_h$ otrzymujemy więc
\[
 \int_\Omega a\nabla r_h\cdot\nabla v\,\mathrm d\lambda_d(x)
 =-\int_\Omega h\nabla\delta y\cdot\nabla v\,\mathrm d\lambda_d(x).
\]
Testując przez $v=r_h$ i wykorzystując koercywność,
otrzymujemy $\|r_h\|_{H_0^1}\le
2\lambda^{-3}\|h\|_\infty^2\|f\|_{H^{-1}}$.
Liniowość $h\mapsto z_h$ wynika z liniowości
prawej strony równania~\eqref{eq:pde32-coeff-linear}
i jednoznaczności rozwiązania. Testując przez
$z_h$, otrzymujemy
$\|z_h\|_V\le\lambda^{-2}\|h\|_\infty\|f\|_{V'}$.
Zbiór współczynników z dodatnim istotnym kresem
dolnym jest otwarty w rzeczywistej przestrzeni
$L^\infty$, gdyż kula o promieniu $\lambda/2$
wokół $a$ jest w nim zawarta.
Oszacowanie reszty dowodzi różniczkowalności
Frécheta w każdym punkcie tego zbioru.
Z reguły łańcuchowej i ciągłego zanurzenia $V$
w $L^2$ wynika także różniczkowalność Frécheta
kosztu $\mathcal F$. Jego pochodna jest reprezentowana
przez funkcję $-\nabla y(a)\cdot\nabla p\in L^1(\Omega)$,
gdyż oba gradienty należą do $L^2$.
Nie wynika stąd jednak istnienie gradientu
w przestrzeni Hilberta $L^2(\Omega)$: iloczyn
dwóch funkcji z $L^2$ nie musi należeć do $L^2$.
\end{proof}

\begin{kontrprzyklad}[Brak identyfikowalności z jednej obserwacji]
Dla $f=0$ i jednorodnych warunków Dirichleta rozwiązaniem
$-\diver(a\nabla y)=0$ jest $y=0$ dla każdego
$a\in L^\infty(\Omega)$ o dodatnim istotnym kresie
dolnym. Istotnie, testowanie równania przez $y$
daje $\lambda\|\nabla y\|_2^2\le0$, gdzie
$\lambda=\operatorname*{ess\,inf}_\Omega a>0$,
a nierówność Poincarégo implikuje $y=0$.
Zatem obserwacja całej
funkcji $y$ nie wyznacza $a$ bez odpowiedniego pobudzenia.
To nie jest wada regularyzacji, lecz strukturalny brak informacji.
\end{kontrprzyklad}

\section{Sterowanie modalne i osiągalność stanów równania ciepła}
\begin{twierdzenie}[Dokładny koszt sterowania pojedynczego modu]
\label{thm:pde32-gramian}
Niech $\lambda>0$, $T>0$ i $a\in\R$. Spośród wszystkich
$h\in L^2(0,T)$, dla których rozwiązanie
$z'(t)+\lambda z(t)=h(t)$, $z(0)=0$, spełnia $z(T)=a$,
minimalną normę ma
\[
 h_*(s)=\frac{a\,e^{-\lambda(T-s)}}{W_\lambda(T)},
 \qquad W_\lambda(T)=\frac{1-e^{-2\lambda T}}{2\lambda}.
\]
Ponadto $\int_0^T|h_*(s)|^2\,\mathrm d\lambda_1(s)
=a^2/W_\lambda(T)$.
\end{twierdzenie}
\begin{proof}
Ponieważ $h\in L^2(0,T)\subset L^1(0,T)$,
funkcja
\[
 z(t)=e^{-\lambda t}\int_0^t e^{\lambda s}h(s)
 \,\mathrm d\lambda_1(s)
\]
jest bezwzględnie ciągła, spełnia równanie
prawie wszędzie i ma wartość $0$ w chwili $0$.
Jednoznaczność wynika z tego, że po pomnożeniu
równania jednorodnego przez $e^{\lambda t}$
pochodna iloczynu znika prawie wszędzie.
Niech $g(s)=e^{-\lambda(T-s)}$. Warunek końcowy
ma postać $(h,g)_{L^2}=a$, a
\[
 \|g\|_{L^2}^2
 =\int_0^T e^{-2\lambda(T-s)}\,\mathrm d\lambda_1(s)
 =W_\lambda(T)>0.
\]
Nierówność Cauchy'ego--Schwarza daje
$|a|^2\le W_\lambda(T)\|h\|_{L^2}^2$.
Funkcja $h_*=ag/W_\lambda(T)$ spełnia warunek
końcowy i osiąga dolne ograniczenie.
Dla każdego innego dopuszczalnego $h$ mamy
$(h-h_*,g)_{L^2}=0$, a stąd
\[
 \|h\|_{L^2}^2
 =\|h_*\|_{L^2}^2+\|h-h_*\|_{L^2}^2.
\]
Dowodzi to jednoznaczności minimalnego sterowania
również dla $a=0$.
\end{proof}

\begin{twierdzenie}[Przestrzeń osiągalnych stanów końcowych]
\label{thm:pde32-reachable}
Niech $\Omega$ będzie niepustym ograniczonym obszarem
Lipschitzowskim, $(e_k)$ bazą ortonormalną $L^2(\Omega)$
złożoną z funkcji własnych laplasjanu Dirichleta,
a $0<\lambda_1\le\lambda_2\le\cdots$ odpowiadającymi
wartościami własnymi liczonymi z krotnościami
(Twierdzenie~\ref{thm:pde18-compact}).
Pracujemy w rzeczywistych przestrzeniach Hilberta.
Rozważmy równanie ciepła $y_t-\Delta y=h$ z zerowym
stanem początkowym i jednorodnym warunkiem Dirichleta,
sterowane na całym obszarze przez $h\in L^2((0,T)\times\Omega)$.
Stan $a=\sum_k a_ke_k\in L^2(\Omega)$ jest osiągalny
w chwili $T>0$ wtedy i tylko wtedy, gdy
\[
 \sum_k\frac{|a_k|^2}{W_{\lambda_k}(T)}<\infty .
\]
Minimalny kwadrat normy sterowania jest równy tej
sumie i jest osiągany przez dokładnie jedno sterowanie.
\end{twierdzenie}
\begin{proof}
Dla $h_k(t)=(h(t),e_k)_{L^2(\Omega)}$ tożsamość
Parsevala i twierdzenie Tonellego dają
\[
 \|h\|_{L^2((0,T)\times\Omega)}^2
 =\sum_k\|h_k\|_{L^2(0,T)}^2.
\]
Z równania słabego testowanego przez $e_k$
wynika $y_k'+\lambda_k y_k=h_k$ w sensie
dystrybucji na $(0,T)$, gdzie $y_k(t)=(y(t),e_k)$.
Współczynniki $y_k$ są bezwzględnie ciągłe:
ich pochodne należą do $L^2(0,T)$, a warunek
$y_k(0)=0$ wynika z ciągłości stanu w $L^2$.
Jeżeli $y(T)=a$, to Twierdzenie~\ref{thm:pde32-gramian}
daje $\|h_k\|_{L^2(0,T)}^2\ge |a_k|^2/W_{\lambda_k}(T)$.
Po zsumowaniu otrzymujemy konieczność warunku
i dolne ograniczenie kosztu.

Odwrotnie, jeśli suma jest skończona, określamy
$h_k^*(s)=a_ke^{-\lambda_k(T-s)}/W_{\lambda_k}(T)$.
Skończone sumy $\sum_{k=1}^N h_k^*(s)e_k(x)$
tworzą ciąg Cauchy'ego w $L^2((0,T)\times\Omega)$,
ponieważ kwadraty norm różnic są ogonami
zbieżnego szeregu kosztów. Ich granica $h^*$
ma współczynniki $h_k^*$; istnieje dla niej
słabe rozwiązanie paraboliczne.
Każdy współczynnik końcowy tego rozwiązania
jest równy $a_k$ z twierdzenia o pojedynczym
modzie. Zupełność bazy daje $y(T)=a$ w $L^2$.
Norma $h^*$ osiąga dolne ograniczenie.
Jednoznaczność wynika z tożsamości
\[
 \|h\|_{L^2((0,T)\times\Omega)}^2
 =\|h^*\|_{L^2((0,T)\times\Omega)}^2+
 \|h-h^*\|_{L^2((0,T)\times\Omega)}^2
\]
dla każdego sterowania osiągającego $a$,
otrzymanej przez zsumowanie tożsamości
dla poszczególnych modów.
\end{proof}
\begin{uwaga}[Zasięg twierdzenia o sterowalności]
Ponieważ $\lambda_k\ge\lambda_1>0$, mamy
\[
 2\lambda_k\le W_{\lambda_k}(T)^{-1}
 \le\frac{2\lambda_k}{1-e^{-2\lambda_1 T}}.
\]
Warunek osiągalności jest więc równoważny
$\sum_k\lambda_k|a_k|^2<\infty$.
Z charakterystyki dziedziny formy Dirichleta
w dowodzie Twierdzenia~\ref{thm:pde18-minmax}
wynika, że zbiór stanów osiągalnych jest dokładnie
przestrzenią $H_0^1(\Omega)$: dla $a$ z tej
przestrzeni mamy
$\|\nabla a\|_2^2=\sum_k\lambda_k|a_k|^2$,
a skończoność szeregu implikuje przynależność do
$H_0^1$ przez zbieżność sum częściowych w normie
formy. W szczególności minimalny kwadrat normy
sterowania, oznaczony przez $M_T(a)$, spełnia
\[
 2\|\nabla a\|_2^2\le M_T(a)
 \le\frac{2}{1-e^{-2\lambda_1T}}\|\nabla a\|_2^2 .
\]
Każdy skończony układ modów jest osiągalny,
a zbiór stanów osiągalnych jest gęsty w $L^2$.
Nie jest on jednak całym $L^2$: wybierając
podciąg $\lambda_{k_j}\ge j^2$ i współczynniki
$a_{k_j}=1/j$, a pozostałe równe zero, dostajemy
$a\in L^2$ oraz $\sum_k\lambda_k|a_k|^2=\infty$.
Dowodzi to sterowalności przybliżonej, ale nie
dokładnej w $L^2$ dla opisanego operatora sterowania.
Nie rozstrzyga to sterowalności przy sterowaniu
ograniczonym do właściwego podobszaru.
\end{uwaga}

\begin{wniosek}[Sterowalność do zera przy sterowaniu na całym obszarze]
\label{cor:pde32-null-control}
Przy założeniach twierdzenia~\ref{thm:pde32-reachable}
dla każdego stanu początkowego $u_0\in L^2(\Omega)$
i każdego $T>0$ istnieje sterowanie
$h\in L^2((0,T)\times\Omega)$ prowadzące do $y(T)=0$.
Istnieje dokładnie jedno takie sterowanie o najmniejszej
normie i spełnia ono
\[
 \|h_*\|_{L^2((0,T)\times\Omega)}^2
 \le T^{-1}\|u_0\|_2^2.
\]
\end{wniosek}
\begin{proof}
Niech $u_{0,k}=(u_0,e_k)_{L^2(\Omega)}$. Z liniowości
równania stan końcowy jest sumą ewolucji swobodnej
$e^{T\Delta_D}u_0$ i stanu osiąganego z zera
przez sterowanie $h$. Aby uzyskać $y(T)=0$,
wystarczy więc osiągnąć stan o współczynnikach
$a_k=-e^{-T\lambda_k}u_{0,k}$ z zerowego stanu
początkowego. Koszt modalny spełnia
\[
 \frac{|a_k|^2}{W_{\lambda_k}(T)}
 =\frac{2\lambda_k}{e^{2T\lambda_k}-1}|u_{0,k}|^2
 \le\frac{|u_{0,k}|^2}{T},
\]
gdzie użyliśmy nierówności $e^r-1\ge r$ dla $r\ge0$.
Sumowanie i tożsamość Parsevala dają skończony
koszt nie większy niż $T^{-1}\|u_0\|_2^2$.
Istnienie i jednoznaczność sterowania o najmniejszej
normie wynikają teraz z twierdzenia
~\ref{thm:pde32-reachable}. Sterowalność do zera
nie przeczy brakowi dokładnej sterowalności w $L^2$:
swobodna ewolucja cieplna wygładza stan początkowy,
a do zera trzeba skompensować jedynie ten wygładzony
stan.
\end{proof}

\section{Podsumowanie metod i ograniczeń}
Regularyzacja stabilizuje problem odwrotny, lecz nie przywraca
informacji utraconej przez operator obserwacji.
W rozważanym kwadratowym zadaniu sterowania optymalnego
istnienie wynika z koercywności, słabej dolnej
półciągłości kosztu i słabej domkniętości zbioru
dopuszczalnego, a jednoznaczność ze ścisłej wypukłości.
Równanie sprzężone pozwala obliczyć pochodną
kosztu bez rozwiązywania osobnego PDE dla każdego kierunku.
Ogólna teoria identyfikacji współczynników z danych brzegowych
(np. problem Calderóna), nierówności obserwowalności,
sterowalność przy sterowaniu na właściwym podobszarze
oraz problemy optymalizacji nieliniowych układów PDE
wymagają dodatkowych metod. Wynik o osiągalności
stanów końcowych dotyczy tutaj równania ciepła
ze sterowaniem rozłożonym na całym obszarze,
a nie ogólnego zagadnienia sterowalności.


\appendix
